% =============================================================================
% §2 Phenomenon & measurement (M1 / R1+R2) — registry collapse; WB/WPP smoothing
% failure; union composition as state variable; data-limitations (measurement).
% =============================================================================
\section{The Phenomenon: Measurement Against Primary Sources}
\label{sec:phenomenon}

\subsection{The collapse in national vital registries}
\label{sec:registries}

Table~\ref{tab:anchors} reports the corrected national anchors. Every figure
comes from the national vital-registration system or statistical institute of
the country concerned; each series in the project repository carries per-row
source and methodology flags.

\begin{table}[t]
\centering
\caption{National TFR anchors, collapse countries and comparator. The CELADE
column reports the corresponding CELADE 2024 medium-variant estimate where a
pinned comparator exists; the systematic register--estimate gap is discussed
in the text.}
\label{tab:anchors}
\footnotesize
\begin{tabular}{@{}l l l c p{4.9cm}@{}}
\toprule
Country & Register series & Level and change & CELADE & Notes \\
\midrule
Costa Rica & INEC & $1.83 \rightarrow 1.12$ (2010--24) & 1.32 &
  register total (includes the 19.3\% of 2024 births to foreign-born mothers);
  denominator INEC-CCP post-Censo-2022, level $\pm$0.03; 2024 from INEC
  \emph{Indicadores demogr\'aficos 2024} \\
Colombia & DANE EEVV & $1.7 \rightarrow 1.1$ (2015--24) & --- &
  \emph{not} the circulating ``$2.0 \rightarrow 1.06$'', which appears in no
  DANE source \\
Chile & INE/DEIS & $1.81 \rightarrow 1.03$ (2010--24) & 1.14 &
  official TGF publication gap 2010--21 filled by births-implied TFR anchored
  to registry years \\
Argentina & DEIS/INDEC & $\approx 1.15$--$1.19$ (2024) & --- &
  no official period TGF after 2021; INDEC-anchored reconstruction, labeled
  throughout \\
Mexico (comp.) & CONAPO / ENADID & 1.89 (modeled) / 1.60 (survey, 2023) & --- &
  slow-decline comparator; the circulating ``1.55'' is not a national figure \\
\bottomrule
\end{tabular}
\end{table}

Three measurement points deserve emphasis. First, the decline is behavioral,
not a denominator artifact: in every collapse country the \emph{raw annual
birth count} fell between 18 and 47 percent while the female population aged
15--49 grew (Chile: 250{,}643 births in 2010 to 154{,}441 in 2024, $-38$
percent; Figure~\ref{fig:forensics}). No plausible revision to population
denominators reverses that.
Second, several widely circulated figures are wrong or unofficial: Colombia's
``2.0 to 1.06'' appears in no DANE source (the genuine collapse is roughly 35
percent in TFR, not 47); Argentina has published no official period TGF since
2021, so any post-2021 level is a reconstruction and should be labeled as such.
Third, the smoothed international series miss the event.
Figure~\ref{fig:smoothing} contrasts national registry series with the World
Bank/WPP series: the model-smoothed lines glide while the registries collapse,
with end-of-sample gaps up to roughly half a child per woman. The same pattern
holds for the regional projections: CELADE 2024 medium-variant estimates sit
\emph{above} the register totals by 0.20 in Costa Rica (1.32 vs.\ 1.12) and
0.11 in Chile (1.14 vs.\ 1.03) --- an optimism gap, not a register deficiency.
The Costa Rican 1.12 is the register \emph{total}: births are a near-complete
civil registry (and include the 19.3 percent of 2024 births to foreign-born
mothers --- immigrant childbearing works \emph{against} the collapse, so the
native-born decline is sharper still); the only revisable component is the
post-Censo-2022 population denominator, worth about $\pm 0.03$ on the level
and nothing on the trend. Analyses built on the international or projected
series would date the collapse late, understate its speed, and miss its
acceleration entirely.

\begin{figure}[t]
\centering
\begin{subfigure}[b]{0.48\textwidth}
  \includegraphics[width=\textwidth]{fig_tfr_registry_wb_COL.png}
  \caption{Colombia}
\end{subfigure}
\hfill
\begin{subfigure}[b]{0.48\textwidth}
  \includegraphics[width=\textwidth]{fig_tfr_registry_wb_CRI.png}
  \caption{Costa Rica}
\end{subfigure}
\caption{National registry TFR versus World Bank/WPP model-smoothed series.
The international series erase the acceleration that defines the phenomenon.}
\label{fig:smoothing}
\end{figure}

\begin{figure}[t]
\centering
\begin{subfigure}[b]{0.48\textwidth}
  \includegraphics[width=\textwidth]{fig_forensic_panel_COL.png}
  \caption{Colombia}
\end{subfigure}
\hfill
\begin{subfigure}[b]{0.48\textwidth}
  \includegraphics[width=\textwidth]{fig_forensic_panel_CRI.png}
  \caption{Costa Rica}
\end{subfigure}
\caption{Registered births and period TFR. Births-implied TFR (anchored to a
registry year) tracks the published national series: the collapse is in the
numerator, not the denominator.}
\label{fig:forensics}
\end{figure}

\subsection{The state variable: union composition}
\label{sec:composition}

What proximately drives births down? We construct annual union-composition
series for women aged 20--39, in four five-year age bands, distinguishing
\emph{married} from \emph{cohabiting} (consensual union) --- for Costa Rica from
ENAHO via INEC's REDATAM tabulations (2010--2024), and for Colombia from DANE
GEIH person-level microdata (2007--2024; eighteen annual pooled files, with the
estado-civil coding verified to be identical across the 2021--22 survey
redesign, so the downturn is not a recode).

The composition series reveal the phenomenon's structure
(Figure~\ref{fig:coupling}). In Colombia, \emph{total} union prevalence sits on
a flat plateau near 59--60 percent from 2008 to 2020 (declining moderately
thereafter), while the \emph{married} share halves --- from 20.7 to 11.9
percent --- with cohabitation substituting almost one-for-one. Costa Rica shows
the same signature: the equal-band-mean married share falls from 0.324 (2010)
to 0.166 (2024). Partnership as such is not disappearing; its \emph{composition}
is transforming, and because fertility differs across union types, composition
is the natural state variable for the fertility consequences. Throughout the
paper, the state variable is the three-way composition
$\{\Sng,\Coh,\Mar\}$ (single, cohabiting, married) --- never a binary
partnered/unpartnered.

\begin{figure}[t]
\centering
\begin{subfigure}[b]{0.48\textwidth}
  \includegraphics[width=\textwidth]{fig_coupling_vs_tfr_COL.png}
  \caption{Colombia (GEIH microdata, 2007--2024)}
\end{subfigure}
\hfill
\begin{subfigure}[b]{0.48\textwidth}
  \includegraphics[width=\textwidth]{fig_coupling_vs_tfr_CRI.png}
  \caption{Costa Rica (ENAHO/REDATAM, 2010--2024)}
\end{subfigure}
\caption{Union composition of women 20--39 (married vs.\ cohabiting) against the
national TFR (comparison overlay only). Marriage collapses on a roughly flat
total-union base; cohabitation substitutes.}
\label{fig:coupling}
\end{figure}

\subsection{Data limitations: measurement}
\label{sec:datalim-measurement}

We flag the measurement limits here rather than in footnotes.
(i)~\emph{Colombia 2007.} The earliest GEIH year is a documented
frame/questionnaire outlier (total union 47.5 percent versus the 2008--2020
plateau near 59--60); all Colombian analysis treats 2008 as the first usable
composition year. (ii)~\emph{Argentina.} No official period TGF exists after
2021; our reconstruction is INDEC-anchored and labeled throughout. The
much-discussed INDEC credibility break of 2007--15 does \emph{not} contaminate
fertility: vital statistics are produced by DEIS, a separate institution.
(iii)~\emph{Chile.} The official TGF series has a 2010--21 publication gap,
filled here by births-implied TFR anchored to registry years.
(iv)~\emph{Pandemic-era collection.} Colombian household-survey collection in
2020--21 switched modes; we flag 2019--22 composition readings as potentially
affected and return to this in Sections~\ref{sec:model}
and~\ref{sec:limitations}. (v)~\emph{Tempo.} Registry TFRs are period measures;
where levels are anchored we also report a tempo-adjusted sensitivity in the
spirit of \citet{bongaarts1998} (a $+0.10$ TFR adjustment changes the Mexican
comparator's implied population by roughly 2 percent); none of the paper's
claims rests on the distinction.
